\section{Дослідження принципів реалізації криптографічних механізмів}

\subsection{Чому при розробці прикладного програмного забезпечення не рекомендується самостійно реалізовувати стандартні криптографічні алгоритми?}

% Відповідь

\subsection{Які основні групи криптографічних примітивів та допоміжних функцій зазвичай надають сучасні криптографічні бібліотеки?}

% Відповідь

\subsection{У чому полягає різниця між криптографічним примітивом та криптографічним протоколом? Навести приклади.}

% Відповідь


\subsection{Що таке криптографічно стійкий генератор псевдовипадкових чисел (CSPRNG) та чому звичайні генератори псевдовипадкових чисел не повинні використовуватися для генерації криптографічних ключів?}

% Відповідь


\subsection{У чому полягає різниця між низькорівневими (low-level) та високорівневими (high-level) криптографічними API? Які ризики може створювати неправильне використання низькорівневого API?}

% Відповідь


\subsection{Що означає реалізація криптографічної операції з постійним часом виконання (constant-time implementation) та від яких атак вона повинна захищати?}

% Відповідь


\subsection{Що таке FIPS 140-3 та в яких випадках наявність відповідної сертифікації криптографічного модуля може бути важливою?}

% Відповідь


\subsection{Чому для реалізації RSA та деяких інших асиметричних криптосистем необхідна арифметика багаторозрядних чисел? Які вимоги, крім швидкодії, висуваються до її реалізації у криптографічному програмному забезпеченні?}

% Відповідь